\documentclass{article}

\PassOptionsToPackage{numbers,compress}{natbib}
% [preprint] = non-anonymous, arXiv-ready (we write alongside development).
% Switch to [final] for camera-ready, or remove the option for anonymous review.
\usepackage[preprint]{neurips_2026}

\usepackage[utf8]{inputenc}
\usepackage[T1]{fontenc}
\usepackage{hyperref}
\usepackage{url}
\usepackage{booktabs}
\usepackage{amsfonts}
\usepackage{amsmath}
\usepackage{amssymb}
\usepackage{nicefrac}
\usepackage{microtype}
\usepackage{xcolor}
\usepackage{graphicx}
\usepackage{subcaption}
\usepackage{float}
\usepackage[capitalize]{cleveref}   % \cref / \Cref for stable cross-refs

\title{Fine-Tuning YOLO11n for Six-Class Waste Detection on Embedded Hardware}

\author{%
  Moussa Ouallaf \\
  \texttt{} \\
}

\begin{document}
\maketitle

% ============================================================================
% The SCRIBE owns the regions between the "% >>> SCRIBE:<zone>" and
% "% <<< SCRIBE:<zone>" markers. Add findings BETWEEN the matching markers.
% Humans: feel free to edit anywhere; the scribe only touches marked zones.
% ============================================================================

\begin{abstract}
% >>> SCRIBE:abstract
% Assembled near the end, once the story is known. Keep the headline number here.
\textit{Abstract to be written once results stabilize.}
% <<< SCRIBE:abstract
\end{abstract}

\section{Introduction}
% >>> SCRIBE:intro
% Problem, why it matters, why it is hard, what we do. End with contributions.
\textit{Introduction in progress.}
% <<< SCRIBE:intro

\paragraph{Contributions.}
\begin{itemize}
% >>> SCRIBE:contributions
  \item \textit{Contributions accrue here as the work lands.}
% <<< SCRIBE:contributions
\end{itemize}

\section{Related Work}
% >>> SCRIBE:related
\textit{Related work accrues here; cite at first mention.}
% <<< SCRIBE:related

\section{Method}
% >>> SCRIBE:method
% finding 1781083278-498726312 (base model decision)
\paragraph{Base detector.}
We fine-tune YOLO11 nano from COCO-pretrained weights~\citep{jocher2024yolo11,lin2014microsoft}.
YOLO11 is the actively maintained Ultralytics mainline.
It outperforms YOLOv8 at fewer parameters~\citep{jocher2023yolov8}.
We rejected the newer YOLO26 for tooling maturity.
YOLO11 small remains the fallback if validation accuracy disappoints.
The deployment target is TensorRT FP16 on a Jetson Orin Nano~\citep{nvidia_tensorrt}.
We expect $\approx 21$\,ms per frame ($\approx 47$\,FPS) at that precision.
This covers three ESP32-CAM streams at 10--15\,fps each.
The detection task itself is deliberately easy.
Each frame shows one object on a white background from a fixed camera.

% finding 1781083278-2727718698 (training data decision)
\paragraph{Training data.}
We train on public waste datasets only and collect no data of our own.
Source labels are remapped to six material classes (plastic, paper, cardboard, metal, glass, organic).
Sources include TACO~\citep{proenca2020taco}, TrashNet~\citep{yang2016trashnet}, WaDaBa~\citep{bobulski2018wadaba}, Kaggle garbage-classification sets and the Drinking Waste Classification set~\citep{serezhkin2020drinking}.

\paragraph{Object presentation and the rest stream.}
Deployment presents one object at a time on a white background under LED lighting.
This staging maximizes transfer from the clean backgrounds of the public sets.
``Rest'' is not a trained class.
Control logic diverts any detection below a confidence threshold to a rest bin.
A trained catch-all class would degrade the six material classes.
% <<< SCRIBE:method

\section{Experimental Setup}
% >>> SCRIBE:setup
% Datasets, baselines, metrics, hyperparameters + how chosen, seeds, compute.
\textit{Experimental setup accrues here.}
% <<< SCRIBE:setup

\section{Results}
% >>> SCRIBE:results
\textit{Results accrue here; report mean $\pm$ std over seeds, tie each number to evidence.}
% <<< SCRIBE:results

% >>> SCRIBE:results-table
% booktabs tables generated from logged metrics go here.
% <<< SCRIBE:results-table

\section{Ablations}
% >>> SCRIBE:ablations
\textit{Ablations accrue here; one row per toggled component.}
% <<< SCRIBE:ablations

\section{Analysis}
% >>> SCRIBE:analysis
\textit{Error analysis, qualitative examples, sensitivity.}
% <<< SCRIBE:analysis

\section*{Limitations}
% >>> SCRIBE:limitations
% Append each caveat the moment it is discovered. No new experiments here.
\textit{Limitations accrue here.}
% <<< SCRIBE:limitations

\section{Conclusion}
% >>> SCRIBE:conclusion
\textit{Conclusion assembled at the end.}
% <<< SCRIBE:conclusion

\section*{Broader Impact}
% >>> SCRIBE:impact
\textit{Societal impact statement.}
% <<< SCRIBE:impact

\section*{Reproducibility Statement}
% >>> SCRIBE:repro
% Point to where each reproducibility element lives (code, seeds, compute, configs).
\textit{Reproducibility pointers accrue here.}
% <<< SCRIBE:repro

\bibliographystyle{unsrtnat}
\bibliography{refs}

\appendix
\section{Appendix}
% >>> SCRIBE:appendix
% Proofs, full hyperparameter tables, datasheets, model cards.
% <<< SCRIBE:appendix

% NeurIPS paper checklist (standalone file written by /paper:init).
% Uncomment for submission:
% \input{checklist}

\end{document}
